本文研究固定随机特征 Galerkin 不可压流求解器的数学适配与底层实现，区分试验空间、时间策略、单次算子成本及运行时开销对性能的贡献。对无反射约束的二维有限圆环流动，给出边界相容的随机特征构造、加权秩截断和含独立均值块的弱离散。推导带边界提升的一次变步 SBDF2 桥接及局部一致性，并在同一角模中构造共同相位坐标，证明几何锚点对指定条件数上界的极小极大性质。工程部分明确 FP64 FFT 打包与共享内存融合、warp 探针归约、C 语言径向内核以及准备和编译缓存的作用。\(Re=100\)、\(T=72\) 下，39 个配置在两种缓存政策下的 78 条 GPU 轨迹均满足速度与涡量 \(0.5\%\) 门槛；分段策略减少真实非线性求值 \(11.11\%\) 和 \(9.36\%\)。融合 FFT 缩短短程推进时间 \(9\%\)--\(17\%\)，却使首次完整求解增加 \(23\%\)--\(69\%\)。新增独立算子验证还显示，实场打包对强弱通道比例敏感，需单独检查弱通道误差。Legendre 降秩对照亦能通过精度门槛，最终隔离计时未被接受，因此尚不能建立随机特征相对同等优化传统 CFD 的整体速度优势。本文提供可复核的离散推导、源码级成本解释与有明确范围的正负结果；局部一致性、变换坐标条件数和数值资格均不被解释为完整非线性稳定性或普适性能保证。
